\section{Common estimates for the two transfers}\label{app:common}
Throughout the transfer appendices $d=m-1$, the reference $g$ and centroids
$c_i$ are as in Section~1, and $K=K_m\geq1$ satisfies~\eqref{eq:Hm}.
Set
\[
 a=|c_i-c_j|=\sqrt2\alpha,\qquad
 P_0=\binom m2,\qquad
 \beta=P_0\sqrt{2/\pi}/a,\qquad \sigma=(K+1)^{-1}.
\]
The first moment obeys $w\leq1$ by Parseval, so
$a^2=2w/(m-1)\leq1$.
All integrals in these appendices use $\gamma_d$ unless specified otherwise.

\subsection{The same nearest frame controls the mismatch}
For balanced $f$ choose a nearest regular centroid matrix $c$ and its partition
$g$. Put $\Delta=w-W(f)$, $q=\norm{b(f)-c}_{\HS}^2$,
$\eta=\frac12\sum_i\norm{f_i-g_i}_1$, and
$J=w-\sum_i c_i\cdot b_i(f)$.
Then
\begin{equation}\label{eq:linear-identity}
 2J=\Delta+q\leq(1+K)\Delta.
\end{equation}
If $\Gamma$ is the gap between the two largest regular linear scores, a union
bound on the $P_0$ Gaussian score differences of standard deviation $a$ gives
$\Prob(\Gamma\leq t)\leq\beta t$. With $p_{\rm wrong}=1-f_{\rm winner}$,
$J\geq\E[p_{\rm wrong}\Gamma]$ and $\E p_{\rm wrong}=\eta$.
Layer cake yields
\[
 J\geq\int_0^\infty(\eta-\beta t)_+\dd t
   =\frac{\eta^2}{2\beta},\qquad
 \eta^2\leq L\Delta,\quad L=\beta(1+K).
\]
Thus no second choice of alignment is needed.

\subsection{A saturated positive gap and a damped negative gap}
Let $\phi,\Phi$ be the standard normal density and distribution function, and
\[
 H_\delta(u)=\phi(u)\Phi(u+\delta)-\phi(u+\delta)\Phi(u).
\]
For $\delta>0$, this is positive: $\Phi(u)/\phi(u)$ is strictly increasing,
as is clear from its integral representation
$\int_0^\infty e^{uy-y^2/2}\dd y$.
In reference cell $i$, put $z_{ij}=(v_i-v_j)\cdot x/\sqrt2\geq0$.
For independent normal utilities with label $i$ of mean zero and label $j$
of mean $-\delta$, the winning-probability difference is
\[
 \E\left[\int_M^\infty H_\delta(u)\dd u\right],
\]
where $M$ is the maximum of the other utilities. Lowering their means
increases this difference.

For $G_i=T_rg_i$, subtract the mean of winner $i$.
The other means are nonpositive and
$\delta=\sqrt2\,r z_{ij}/\sqrt{1-r^2}$.
Raise all other means to zero; the event $M\leq0$ has probability
$2^{-(m-2)}$. For $0\leq u,\delta\leq1$,
$\partial_\delta H_\delta(u)\geq\phi(1)\phi(2)$.
Integrate over $0\leq u\leq1$; for larger $\delta$ use monotonicity when the
competitor mean is lowered. Consequently
\begin{equation}\label{eq:plus-gap}
 G_i-G_j\geq c_+\min(rz_{ij},1),\qquad
 c_+=\frac{e^{-5/2}}{2^{m-1}\pi}.
\end{equation}

For $G_i=T_{-r}g_i$, subtract label $j$'s mean instead.
Now $j$ has mean zero, $i$ has mean $-\delta$, and every other mean is at most
$U_x=\sqrt2\,r|x|/\sqrt{1-r^2}$. Raise those means to $U_x$.
Their maximum is at most $U_x$ with probability $2^{-(m-2)}$.
On $U_x\leq u\leq U_x+1$ and $0\leq\delta\leq1$,
\[
 \partial_\delta H_\delta(u)
 =\phi(u)\phi(u+\delta)+(u+\delta)\phi(u+\delta)\Phi(u)
 \geq\phi(U_x+1)\phi(U_x+2).
\]
The same integral argument and monotonicity for $\delta\geq1$ give a gap at
least $c_+e^{-U_x^2-3U_x}\min(\delta,1)$. Since
$3U_x\leq U_x^2+9/4$ and
$2U_x^2=4r^2|x|^2/(1-r^2)\leq|x|^2/32$ for $r\leq1/16$,
\begin{equation}\label{eq:minus-gap}
 G_j-G_i\geq c_-\min(rz_{ij},1)e^{-|x|^2/32},
 \qquad c_-=\frac{e^{-19/4}}{2^{m-1}\pi}.
\end{equation}
Thus the reference winner is the unique largest positive potential and the
unique smallest negative potential, outside null hyperplanes.

\subsection{Weighted moments and uniform tails}
Let
\[
 \psi(x)=e^{|x|^2/8},\qquad
 \omega_r(x)=\min\left(r\min_{i<j}|(v_i-v_j)\cdot x|/\sqrt2,\,1\right).
\]
Use $\widetilde\omega_r^+=\omega_r$ and
$\widetilde\omega_r^-=e^{-|x|^2/32}\omega_r$.
Define the finite constants
\begin{align*}
 B_+&=P_0\left[2^{(d-1)/2}\phi(0)e^{1/4}+2^{d/2-2}\right],\\
 B_-&=P_0\left[\phi(0)e^{5/16}(8/3)^{(d-1)/2}
                  +\tfrac14(8/3)^{d/2}\right].
\end{align*}
For $|h|\leq1$, $\lambda>0$, and $0<r\leq1$,
\begin{equation}\label{eq:Fenchel}
 \int|h|\psi\leq\lambda\int|h|\widetilde\omega_r^\pm
                      +\frac{B_\pm}{\lambda r}.
\end{equation}
Here are the constants explicitly. For the positive case dominate the
positive part by summing over the individual planes. Write $x=ue+z$.
On $|u|\leq1$, replace $\psi$ by
$A(z)=e^{1/8}e^{|z|^2/8}$ and $\phi(u)$ by $\phi(0)$, and integrate
$(A-\lambda r|u|)_+$ over the whole line. Its integral is
$A^2/(\lambda r)$. On $|u|>1$ the saturated plane weight is at least $r$,
and $(A-B)_+\leq A^2/(4B)$. Gaussian exponential moments give $B_+$.
For the negative case use the exact identity
\[
 (\psi-\lambda\widetilde\omega_r^-)_+
 =e^{-|x|^2/32}(e^{5|x|^2/32}-\lambda\omega_r)_+
 \leq(e^{5|x|^2/32}-\lambda\omega_r)_+.
\]
The same calculation, with this auxiliary exponential, gives $B_-$.
The moment controlled on the left of~\eqref{eq:Fenchel} remains $\psi$.

For any relaxed $k$, define the nonnegative losses
\[
 D_+(k)=-\sum_i\ip{k_i-g_i}{T_rg_i},\qquad
 D_-(k)= \sum_i\ip{k_i-g_i}{T_{-r}g_i}.
\]
Each vanishes only for $k=g$ by the strict gaps. Optimizing
\eqref{eq:Fenchel}, using $\sum_i|k_i-g_i|=2p_{\rm wrong}$ and
\eqref{eq:plus-gap} or~\eqref{eq:minus-gap}, gives
\begin{equation}\label{eq:weighted-moments}
 \sum_i\left(\int|k_i-g_i|\psi\right)^2\leq M_\pm D_\pm(k)/r,
 \qquad M_\pm=8B_\pm/c_\pm.
\end{equation}
In the relevant sign put $A_\pm^2=\max(1,m/4)M_\pm$.
If $z=\sqrt{D_\pm(k)/r}$ and $\eta=\frac12\sum_i\norm{k_i-g_i}_1$,
then $\eta\leq A_\pm z$.
For $R\geq2$, let $\eps_R=(1+R)e^{-R^2/8}$.
The pointwise bound
$(1+|x|)\1_{\{|x|>R\}}\leq\eps_R\psi(x)$ implies both
\begin{equation}\label{eq:tails}
 \frac12\sum_i\left|\E[(k_i-g_i)\1_{\{|X|>R\}}]\right|
 \leq\eps_RA_\pm z,\qquad
 \left(\sum_i\left|\E[X(k_i-g_i)\1_{\{|X|>R\}}]\right|^2\right)^{1/2}
 \leq\eps_RA_\pm z.
\end{equation}
These are separate bounds, obtained from Cauchy--Schwarz across labels.

\subsection{The quadratic Mehler error}
The Mehler kernel relative to $\gamma_d\otimes\gamma_d$ is
\[
 M_\rho(x,y)=(1-\rho^2)^{-d/2}
 \exp\!\left(\frac{2\rho x\cdot y-\rho^2(|x|^2+|y|^2)}
                   {2(1-\rho^2)}\right).
\]
For $|\rho|\leq1/16$, writing $\ell=\log M_\rho$ and
$L_{xy}=|x|^2+|y|^2$, direct differentiation gives
$|\ell'|\leq d+L_{xy}$, $|\ell''|\leq2d+2L_{xy}$ and
$M_\rho\leq2^d e^{L_{xy}/34}$.
For example the nondimensional parts of $\ell'$ and $\ell''$ are
\[
 \frac{(1+\rho^2)x\cdot y-\rho L_{xy}}{(1-\rho^2)^2},
 \qquad
 \frac{2\rho(3+\rho^2)x\cdot y-(1+3\rho^2)L_{xy}}{(1-\rho^2)^3};
\]
$|x\cdot y|\leq L_{xy}/2$ gives the stated bounds. The exponent bound follows
from $|\rho|/[2(1+|\rho|)]\leq1/34$.
Put $t_0=13/136$ and
\[
 C_d=2^{d-1}\left[d^2+2d+\frac{2d+2}{2t_0}+\frac1{t_0^2}\right].
\]
Taylor's theorem, $M''=M((\ell')^2+\ell'')$, and the maxima of
$L_{xy}e^{-t_0L_{xy}}$ and $L_{xy}^2e^{-t_0L_{xy}}$, using $e>2$, give
\begin{equation}\label{eq:Mehler}
 |M_\rho-1-\rho x\cdot y|
 \leq C_d\rho^2\psi(x)\psi(y).
\end{equation}
Thus a mixed remainder for deviations $f-g,h-g$ is bounded by
$C_dM_\pm r\sqrt{D_\pm(f)D_\pm(h)}$.

\subsection{Exact mass repair on fixed facet slabs}\label{subsec:repair}
Fix a reference frame, $R\geq3$, and a relaxed $k$ with mismatch $\eta$.
Let $p=\E k$, $U=\1/m$, $u=\frac12\norm{p-U}_1$.
Truncate $k$ to $g$ outside $B_R$, obtaining $k^R$, whose half-$\ell^1$
mass imbalance $v$ satisfies
\[
 v\leq u+\eps_RA_\pm z,\qquad v\leq\eta.
\]
Set
\[
 c_A=\Vol(B_{d-1})\,8^{-(d-1)}(2\pi)^{-d/2}e^{-2}>0.
\]
For donor $i$ and recipient $j$, use the facet disk centered at $v_i+v_j$,
of tangential radius $1/8$, and its donor-side normal slab of width
$\tau=(\eta+2v)/c_A$.
Assume $\eta\leq c_A/24$, so $\tau\leq1/8$.
At the anchor, the \emph{unscaled} $v_i$ and $v_j$ scores tie and exceed
all other scores by one. A displacement of size at most $\sqrt2/8$ reduces
these gaps by at most $1/4$. Distinct unoriented anchors are at least
$\sqrt2$ apart; opposite slabs at one anchor have disjoint interiors.
All oriented supports therefore lie in their correct cells, are mutually
disjoint up to null sets, and are contained in $B_2$.

Their Gaussian mass is at least $c_A\tau$.
Since the total missing donor mass in its reference cell is at most $\eta$,
each slab has available $k_i^R$ mass at least $c_A\tau-\eta=2v$.
Choose any excess-to-deficit flow of total mass $v$, and transfer each pair
mass proportionally to $k_i^R$ on its slab. Each pair amount is at most $v$.
Disjoint supports keep the assignment simplex-valued, and all masses are
repaired exactly. If $v=0$, make no transfers. Call the result $k'$.

Let $b'=b(k')$, $e_0=\norm{b'-b(k)}_{\HS}$, and
\[
 J'=w-\sum_i c_i\cdot b_i',\qquad
 J_{\rm in}=\E\!\left[\1_{B_R}\sum_i k_i(X)
                  (\max_j c_j\cdot X-c_i\cdot X)\right].
\]
The linear cost per unit transferred mass is at most $a\tau\leq\tau$.
Its total is at most $3\eta v/c_A$. Tail removal and the centroid cost of
transferring mass in $B_3$ consequently give
\begin{equation}\label{eq:repair-bounds}
 J'\leq J_{\rm in}+\frac3{c_A}\eta u+C_J\eps_Rz^2,\qquad
 e_0\leq6(u+\eps_RA_\pm z),\qquad C_J=3A_\pm^2/c_A.
\end{equation}
For the first inequality use $v\leq u+\eps_RA_\pm z$ and
$\eta\leq A_\pm z$; for the second use~\eqref{eq:tails} and
$3\sqrt2v$ as the Frobenius centroid cost.

Suppose now that $c$ is a nearest regular centroid matrix for $b(k)$,
and put $q=\norm{b(k)-c}_{\HS}^2$.
Distance to $\cO_m$ is one-Lipschitz. Applying~\eqref{eq:Hm} to balanced $k'$
gives
$2J'\geq(1+1/K)(\sqrt q-e_0)_+^2$.
Young's inequality with parameter $1/(2K)$, also covering $\sqrt q\leq e_0$,
therefore gives
\begin{equation}\label{eq:repair-coercivity}
 q\leq(2-\sigma)J'+(2K+1)e_0^2.
\end{equation}
The linear cost retains the original frame; no new alignment is introduced.

\subsection{Comparison on a fixed ball}
Write $H_i(y)=\E g_i(y+Z)$. Then $\nabla H_i(0)=c_i$ and
$\norm{\nabla^2H_i}_{\rm op}\leq\sqrt2$, by Gaussian differentiation:
in a unit direction the second derivative is
$\E[g_i(y+Z)((u\cdot Z)^2-1)]$, bounded by Cauchy--Schwarz.
Conicality gives $T_{\pm r}g_i(x)=H_i(\pm tx)$,
$t=r/\sqrt{1-r^2}$. On $B_R$ and for $r\leq1/16$,
\[
 \norm{\nabla(T_{\pm r}g_i-1/m\mp r c_i\cdot x)}
 \leq(1+3R)r^2.
\]
The orthogonal reflection swapping vertices $i,j$ fixes their bisector and
all other vertices, so the pair remainder vanishes on the bisector.
Integrating normally from that plane, inside $B_R$, gives the appropriate
signed potential gap at least
$r(c_i-c_j)\cdot x(1-\beta_Rr)$ in reference cell $i$, where
\begin{equation}\label{eq:compact}
 \beta_R=\frac{2(1+3R)}a,\qquad
 rJ_{\rm in}\leq\frac{D_\pm(k)}{1-\beta_Rr}
 \quad\hbox{if }\beta_Rr<1.
\end{equation}

